\section{尝试策略}

\subsection{策略一：上下文裁剪}

对历史对话进行裁剪，只保留与当前轮请求直接相关的内容。该策略可以降低 prefill token 数，从而减少 TTFT，但需要避免删掉影响任务理解的关键信息。

\subsection{策略二：稳定 Prompt 前缀}

将系统提示词、工具说明和固定格式模板放在稳定前缀中，减少每轮请求中前缀内容的变化。稳定前缀有利于 prefix KV cache 命中，适合多轮对话 Agent 的连续请求场景。

\subsection{策略三：会话级 KV Cache 复用}

为同一个会话维护可复用的 KV cache，使新一轮请求只对新增用户输入和必要增量上下文进行 prefill。该策略对降低多轮 TTFT 最直接，但需要管理 cache 生命周期、显存占用和会话淘汰策略。

\subsection{策略四：工具调用后置或异步化}

对于不影响首轮响应的工具调用，尝试后置或异步执行，先让模型返回初始响应或状态提示。该策略可以降低端到端 TTFT，但只适用于工具结果不是首 token 必需条件的场景。

\subsection{策略五：请求调度优化}

通过限制最大 batch 等待时间、区分短上下文和长上下文请求、为交互式请求设置更高优先级等方式降低排队时间。该策略主要缓解高并发下 TTFT 抖动。
多轮 Agent 场景中，不同请求的输入长度和输出长度差异很大。长上下文请求会带来较大的 prefill 开销，如果与短请求混在同一入口直接进入推理服务，可能导致短请求被长请求阻塞，首 token 延迟明显上升。因此，可以在 server 层进行准入控制：限制同时进入 vLLM 的长请求数量，为短请求保留快速通道，并对等待时间过长的请求进行补偿，避免长期饥饿。
该策略不替代 vLLM 内部调度器，而是在请求进入引擎前进行节流和分流。server 层可以根据输入 token 数、历史 TTFT、当前队列长度、backend 状态和请求阶段决定是否立即放行请求。对于交互式、多轮短请求，可以提高优先级；对于长上下文整合请求，可以限制并发数，减少对其他请求的影响。


\subsection{策略六：Tool call cache}
在多轮 Agent 任务中，模型可能反复调用搜索、文档查询、元数据读取、只读检索等工具。如果每次都重新执行相同或高度相似的工具调用，会增加额外等待时间，并拉长整轮 Agent 的端到端耗时。因此，我们尝试在模型发出工具调用和实际执行工具之间加入 Tool Call Cache，用于复用已有工具返回结果。

Tool Call Cache 与 KV Cache 作用层次不同。KV Cache 解决的是模型内部重复上下文带来的 prefill 计算问题；Tool Call Cache 解决的是模型外部重复工具调用带来的等待和 I/O 开销问题。相比直接修改 vLLM 底层，Tool Call Cache 可以在 server 层或 Agent 工具链中实现，风险更低，也更适合作为当前阶段的落地策略。

在具体设计上，我们采用风险分层策略：对于文档检索、只读查询、静态元数据读取等低风险工具，可以加入白名单并设置较长 TTL；对于搜索结果、天气、金融信息等具有时效性的中风险工具，需要设置较短 TTL；对于写文件、删除文件、执行命令、发送请求、提交代码等具有副作用的高风险工具，默认绕过 cache。

cache key 由工具名称、标准化后的工具参数和必要上下文状态共同构成。参数标准化时需要去除请求编号、trace id、时间戳等不影响工具语义的字段，同时保留真正决定工具行为的参数。对于依赖文件状态或工作区状态的工具，还需要加入文件摘要、修改时间或会话状态，避免文件内容变化后仍复用旧结果。

此外，在高并发场景下，相同工具调用可能被多个 session 同时触发。为减少重复执行，引入了并发去重机制：第一个请求真正执行工具，后续相同 cache key 的请求等待同一份执行结果，完成后共同复用。配合 TTL、LRU 淘汰和命中率统计，可以进一步评估该策略对工具调用延迟和端到端耗时的影响。

\subsection{策略七：GDN下system prompt和user prompt的区分}
在多轮 Agent 请求中，system prompt 通常包含任务规则、工具说明和输出格式，具有较强稳定性；user prompt 则随每轮输入变化。二者的复用价值不同：system prompt 更适合作为稳定前缀保留，user prompt 更适合作为增量内容处理。因此，我们尝试在 server 层区分 system prompt 和 user prompt，为后续上下文组织、prefix cache 友好路由和 cache-aware 调度提供信号。

在 GDN 场景下，这一策略需要保持谨慎。由于 GDN 相关机制与标准 Transformer 的 KV Cache 复用路径并不完全一致，不能简单认为识别出 system prompt 后就可以直接复用对应 KV Cache。因此，我们主要在 server 层做轻量识别：根据角色字段、prompt 模板、固定工具说明、公共前缀长度和多轮出现频率，判断哪些片段属于稳定 system prompt，哪些片段属于当前轮 user prompt 或动态上下文。

具体实现上，可以为每个会话维护 prompt 状态记录，包括稳定前缀摘要、system prompt 长度、公共前缀命中次数、最近几轮 user prompt 变化情况和工具结果写回位置。新请求到达时，server 层解析请求结构并更新状态。如果某段前缀在多轮中高度一致，则标记为稳定 system prompt；如果某段内容随轮次频繁变化，则归为 user prompt 或动态上下文。

该策略的作用主要体现在三方面：一是保持 system prompt 的位置和文本稳定，减少对公共前缀的破坏；二是将 system prompt 摘要作为多 backend 路由信号，使相似请求更容易聚集到同一 backend；三是在上下文裁剪时优先保留 system prompt 和长期约束，对 user 历史和工具结果进行预算化管理。需要注意的是，该区分本身不会直接减少计算，真正的 prefill 降低仍依赖 vLLM 底层 cache 机制配合。

\subsection{策略八：多 backend 覆盖度路由}

当系统中存在多个 vLLM backend 时，如果请求随机分发，每个 backend 上的 prefix cache 局部性会被破坏，相似请求可能被分散到不同实例，导致重复 prefill。为此，可以引入 cache-aware 的多 backend 覆盖度路由：新请求到达时，不只判断哪个 backend 空闲，还估计该 backend 近期处理过的请求能否覆盖当前请求。

由于 server 层无法直接读取 vLLM 内部 KV Cache 内容，可以采用轻量近似方法：为每个 backend 维护近期请求摘要，包括 system prompt 摘要、prefix block 摘要、工具结果摘要、session 标识、请求阶段、近期 TTFT 和队列状态。新请求到达后，根据其 system prompt、前缀内容、工具结果和 session 信息与各 backend 的近期记录进行匹配，计算覆盖度分数。覆盖度越高，说明该 backend 更可能保留相关 cache，适合优先路由。

同时，覆盖度路由不能只追求 cache 命中。如果某个 backend 覆盖度高但当前队列较长，继续把请求分发过去反而会增加端到端延迟。因此需要加入延迟 gate，根据当前排队长度、近期 TTFT、显存压力和请求负载对 backend 进行降权或过滤。最终路由决策应在 cache 覆盖收益和延迟成本之间做权衡。

\subsection{策略九：超参数搜索与配置选择}

除了面向 Agent 结构的策略外，我们也尝试通过超参数实验寻找更适合当前硬件和负载的 vLLM 配置。核心参数包括 TP、PP、DP、max model length、GPU memory utilization、并发请求数和长请求比例等。

TP、PP、DP 的不同组合会影响模型能否稳定启动、显存是否足够、通信开销大小、吞吐和首 token 延迟。max model length 决定系统能够承载的上下文长度，也直接影响 KV Cache 空间需求；GPU memory utilization 影响可分配给 KV Cache 的显存比例；并发请求数和长请求比例则会改变排队时间和整体吞吐。

通过自动化扫参和多轮负载测试，可以将手工试错变成可复现的实验过程。每个实验点记录配置、日志、错误信息和性能指标，从而分析不同参数组合对 TTFT、端到端延迟、吞吐、显存占用和上下文裁剪比例的影响。该策略本身不改变模型或 Agent 逻辑，但可以为后续 cache、路由和调度策略提供稳定 baseline。
